\documentclass[10pt,a4paper]{amsart}
\usepackage[margin=19mm]{geometry}
\usepackage{amsmath,amssymb,mathtools}
\usepackage[scr=rsfs,cal=euler]{mathalfa}
\usepackage{xfrac}
\usepackage[unicode,hidelinks,bookmarks=false]{hyperref}
\newcommand{\ie}{i.e.\ }
\newcommand{\cf}{cf.\ }
\newcommand{\resp}{respectively\ }
\newcommand{\Subsection}{\subsection}
\providecommand{\nobreakdash}{\nobreak\hbox{-}}
\setlength{\emergencystretch}{2em}
\multlinegap0pt
\usepackage{amssymb}
\usepackage[shortlabels]{enumitem}
\newtheorem{thm}{Theorem}
\title{Arakeljan's Theorem with bounded approximants}
\author{Spyros Pasias}
\date{Source: arXiv:2304.03334v7 (CC BY 4.0)}
\begin{document}
\maketitle

\noindent\emph{Source and licence.} Arakeljan's Theorem with bounded approximants. Version arXiv:2304.03334v7 (CC BY 4.0), distributed by arXiv under the Creative Commons Attribution 4.0 International licence, http://creativecommons.org/licenses/by/4.0/, as recorded in the arXiv licence metadata for this exact version and shown on the abstract page https://arxiv.org/abs/2304.03334v7. Full licence notice: \texttt{licenses/arxiv-2304.03334-CC-BY-4.0.txt}.\par\smallskip

\noindent\textit{Editorial scope.} Counted unit: the article's thm thm (source0.tex lines 208--213). The article's complete original proof of it is delivered with it (source0.tex lines 216--236, one \texttt{proof} environment of 21 lines). No mathematics is written, repaired or re-derived: the statement and the proof are the article's own lines, in the article's own notation, and no retained line is substituted or re-flowed.  No further source-local context is needed: every cross-reference of every retained line resolves inside the two delivered spans.  Nothing was omitted from any retained span, so the package carries no omission record.  No retained line contains a citation, so the wrapper carries no bibliography and no bibliography entry.  No retained line contains a figure environment, an image-inclusion command or drawing code, so no image is delivered and the package declares no asset dependency.  Editorial formatting only, in addition: an AMS \texttt{amsart} preamble that restates the article's own declarations and macros used by the retained lines, namely the environments \texttt{thm} on separate counters, together with the standard packages those lines need.  The retained environments print in this document as Theorem 1 in order of appearance, whereas the article prints the same items with the numbering of its own full document.  The complete native source of the article as distributed in the arXiv e-print for this exact version is bundled unchanged as \texttt{sources/139793/source0.tex}.
\begin{thm}[Main Theorem]
Let $G\subset \mathbb{C}$ be an arbitrary domain and suppose $F$ is a relatively closed subset of $G$, and $C$ is a closed subset in $G$ such that \(F\) and \(C\) are \(G-\)hole independent. In order that every function $f\in A(F)$ can be approximated uniformly on $F$ by functions $g\in Hol(G)$, each of which is bounded on $C$, it is necessary and sufficient that: \begin{enumerate}
\item $F$ is an Arakeljan set.
\item There exists a closed Arakeljan set $C_1\not=G$, so that $C\subset C_1$ and $F\cap C_1$ is a bounded set.
\end{enumerate}
\end{thm}

\begin{proof}[Main Theorem]


\textbf{(Necessity.)} Assume that any $f\in A(F)$ can be approximated uniformly by functions $g\in Hol(G)$, each of which is bounded on $C$. By Arakeljan's Theorem it follows $F$ is an Arakeljan set. Clearly the function $\phi(z)=z$ belongs to the class $Hol(G)$, thus there exists a function $f\in Hol(G)$ such that: \begin{equation}|z-f(z)|<1,\   \ z\in F  
\end{equation}and $f$ is also bounded on C.
 By {lemma 1} there exists a closed Arakeljan set $C_1\subset \mathbb{C}$ such that that $C\subset C_1$ and $f(z)$ is bounded on $C_1$. By $(1)$ clearly $F\cap C_1$ is bounded. The proof of necessity is complete.
 
 \noindent \textbf{(Sufficiency.)} We have Arakeljan sets $F$ and $C_1$, $C\subset C_1$, such that $F\cap C_1$ is a bounded set. Let $K$ be a compact set in $G$ such that $\partial K$ is the union of jordan curves and so large that:
 \begin{equation} 
  F\cap C_1\subset K
   \end{equation}
  \begin{equation}
  (\partial K)\setminus C_1\not=\emptyset
   \end{equation}

\noindent Note that $(3)$ is possible since $C_1\not=G$ is a closed set without $G-$holes.\\
 The set $C_1\setminus K$ is also without $G-$holes, because otherwise there would exist a $G-$hole $g$ of $C_1\setminus K$. This leads to a contradiction since we must have $\partial g\subset C_1$. However, since $C_1$ has no $G-$holes it follows that $g=K$ and this contradicts $(3)$.\\
 The set $C_1\setminus K$ is also an Arakeljan set. Otherwise, there would exist a compact set $L\subset G$ so that the set $H\equiv \lbrace h: h \text{ is a }G-\text{hole of }L\cup(C_1\setminus K)\rbrace$ is either unbounded or $\partial H\cap \partial G\not=\emptyset$. But this is a contradiction because then the set of $G-$holes of the set $C_1\setminus K$ union the compact set $L\cup K$ would also have to either be unbounded or accumulate to a point in $\partial G$. Indeed, the claim follows, since $(L\cup K)\cup (C_1\setminus K)=C_1\cup K$. This contradicts that $C_1$ is an Arakeljan set, as such, by the contradiction we conclude $C_1\setminus K$ is an Arakeljan set.\\
 Now the Arakeljan sets $F$ and $C_1\setminus K$ are disjoint by $(2)$, therefore by the previous lemma we have that $F\cup (C_1\setminus K)$ is an Arakeljan set. \\
 Let $\phi(z)\in A(F)$ be any function. Define a function $h(z)$ by $h(z)=\phi(z)$ on $F$ and $h(z)=0$ on $C_1\setminus K$. Clearly by $(2)$ and the lemmas above it is evident that $f\in A(F\cup (C_1\setminus K))$. Therefore by Arakeljan's theorem $f$ is uniformly approximable by functions in $A(G)$. Therefore for any $\epsilon>0$, there exists a function $f\in A(G)$ such that $|h(z)-f(z)|<\epsilon$ for any $z\in F\cup (C_1\setminus K)$. Hence we have that $|\phi(z)-f(z)|<\epsilon$ on $F$ and $|f(z)|<\epsilon$ on $C_1\setminus K$. The function $f(z)$ is bounded on $C_1\setminus K$, moreover $f(z)$ is bounded on the compact set $C_1\cap K$ thus, $f(z)$ uniformly approximates $\phi(z)$ on $F$ and is bounded on $C$. The proof is complete.
\end{proof}

\end{document}
